\chapter{Introduction to \LaTeX}
\LaTeX{} is a text formatting system created by Leslie Lamport in the early 1980s\cite{latex}. The program is based on the \TeX{} text formatting system created by Donald Knuth in 1978\cite{tex}. With 
\LaTeX{} you author your document by editing a text input file using a
program such as EMACS, |vim|, or another editor. You then give
this input file to \LaTeX{} (or, more accurately, to a program such as
|pdflatex| or |xelatex|). \LaTeX{} then transforms your
input file(s) into an Adobe Portable Document Format (PDF) file.

Although most documents at NPS are prepared with
Microsoft\textregistered{} Word, \LaTeX{} is widely used outside of
NPS in the sciences to create conference papers, journal articles, and
even full-length books. \LaTeX{} is especially popular in computer
science. With the NPS template you can use \LaTeX{} to produce an NPS
thesis that is consistent with NPS formatting requirements. Because
\LaTeX{} provides for automated formatting, automatic updating of
references, and the ability to directly embed experimental results,
many students who write technical documents at NPS that use \LaTeX{}
find that they save time---even when they take into account the time
that they spend learning to use \LaTeX{} in the first place!

\section{Reasons to use \LaTeX}
There are many reasons to use \LaTeX{} for preparing
a technical document:

\begin{itemize}
\item Formatting for paragraphs, quotations, lists, tables, and other
  elements is performed automatically. You can easily make changes to
  formatting and have them reflected throughout your document.  The
  numbers used for numbered sections are automatically updated when
  new material is added or removed. As a result, formatting is more
  consistent with \LaTeX{} than with Microsoft Word.

\item References within your document to numbered chapters,
  appendices, sections, figures, tables, equations and references are
  updated automatically each time a |.pdf| file is produced, assuring
  that they are  correct.

\item The Bib\TeX{} system produces consistent citations and
  bibliography. References are pulled from a bibliographic database
  that is separately maintained and can be shared between many documents. Records for the bibliographic
  database can be downloaded from many online services, helping to
  assure that they are consistent. The citation format is maintained 
  separately from the citation contents, making it easy to change
  citation styles when submitting to different conferences or journals.  
\item \LaTeX{} allows you to directly include other files at the time
  that the |.pdf| file is created. This makes it easy to automatically
  incorporate the results of experiments in tabular or graphical form,
  without having to manually copy results.  Source code for programs
  can be included with pretty syntax from the original files, and line
  numbers can be automatically displayed as desired. 
\item |pdflatex|, the version of \LaTeX{} we recommend at NPS, also
  allows you to embed other files as attachments within your |.pdf| file. This makes it
  easy to preserve  experimental data, spreadsheets, or other
  information in the final file that is distributed to sponsors and archived. 
\item Because the \LaTeX{} input file is plain ASCII, you can store
  your document using a revision control system such as
  Subversion (SVN)\cite{subversion}. This allows multiple people to work on
  the same document at the same time; Subversion automatically merges
  the changes together.  If you save your thesis work daily, it is possible
  retrieve previous revisions of your thesis and undo changes or
  mistakes---even many months after the fact.
\item \LaTeX{} is free software and runs on PCs, Macs, and Unix
  systems. This means that you can produce your documents on
  practically every computer you have, without having to purchase
  anything else.
\end{itemize}

NPS has developed this template for preparing NPS reports and theses. If
you are using \LaTeX{} at NPS, there are many reasons for using the
template:

\begin{itemize}
\item The template has been in use since 2007. Numerous NPS
  faculty members have worked with the NPS Thesis Processor to assure
  that the documents produced are acceptable for student and faculty
  use.
\item The template supports the creation of unclassified, For Official Use Only (FOUO), and
  classified documents. In particular, the template supports
  classification labels for paragraphs, captions, and references.
\item The cover page, signature page, and Standard Form 298 are automatically generated.
\item Combined with the Subversion, the template makes it possible for students
  and their advisors to collaborate on a document without the need to pass files back and forth. 
\end{itemize}


\section{The Purpose of this Document}

This document does not duplicate the depth of information
available elsehwere on \LaTeX{},
but does provide you with the minimum amount of information required to use \LaTeX{}
to produce a master's thesis or technical report at NPS.

If you wish to learn more about \LaTeX, there are many good reference
books and online tutorials for \LaTeX{}.

We recommend these online resources:
\begin{enumerate}
\item \url{http://en.wikibooks.org/wiki/LaTeX}, a wikibook that
  discusses many aspects of \LaTeX.
\item \url{http://www.ctan.org/}, the Comprehensive \TeX{} Archive
  Network, a collection of modules and documentation for extending \LaTeX{}.
\item \url{http://en.wikipedia.org/wiki/Comparison_of_TeX_editors}, a
  web page on Wikipedia that discusses different programs available for editing \LaTeX{} input files.
\item \url{http://tex.stackexchange.com/}, a question-and-answer website for \TeX{} and \LaTeX{} problems.
\end{enumerate}

We also recommend these books. They are expensive, but worth if it if
you wish to become a \LaTeX{} master.
\begin{enumerate}
\item Guide to \LaTeX{} (4th Edition), by Patrick W. Daly. This is the
  comprehensive \LaTeX{} reference which will provide you with an
  astounding amount of \LaTeX-related information. Do not read
  Lamport's original book, as it is quite out of date.
\item The \LaTeX{} Companion (Tools and Techniques for Computer
  Typsetting), by Frank Mittelbach, Michel Gossens, Johannes Braams,
  David Carlisle and Chris Rowley.
\item The \LaTeX{} Graphics Companion, by Michael Gossens, Frank
  Mittelbach, Sebastian Rahtz, and Denis Roegel.
\end{enumerate}

\section{Installation}
Before you can use \LaTeX{}, you will need to install two critical
pieces of software:
\begin{enumerate}
\item The \LaTeX{} system itself.
\item A program for editing the \verb+.tex+ input files.
\end{enumerate}

Here once again there are many options. For both \LaTeX{}
and text editors there are both free  and commercial
distributions. This document makes specific recommendations that were
known to work as of the document's date of publication. You are free
to explore on your own as well.

\subsection{Installation on MacOS 10.5 and above}

The easiest way to get \LaTeX{} operational on a Mac is
 to download an installer for the most recent distribution
 from the \TeX Users's Group (TUG) at \url{http://www.tug.org/mactex}. 

TUG's Mac\TeX{} distribution will install \LaTeX{} in the
\verb+/usr/texbin/pdflatex+ directory and will update your startup files to
include this directory in your path. If you chose this strategy, be
sure to click ``Customize'' in the installer and select  all of
the optional packages for installation.

You can also install the program from sources using the MacPorts or Fink installer system, but this is not recommended. 

If you are using \LaTeX{} on a Mac, you should consider downloading
and using \emph{LaTeXiT}, an open-source program that allows you to type
\LaTeX{} math and create PDFs for embedding in other applications. The
created PDFs have the source \LaTeX{} code embedded, so you can re-edit
them in LaTeXiT using the Mac ``Services'' feature.  You may also consider the
commercial program \emph{Latexian}, which allows you to type a
\LaTeX{} document in one window
and see the PDF update in another window as you type.

\subsection{Installing on Linux}
For most Linux systems a complete \LaTeX{} distribution 
can be downloaded as part of the |texlive-latex3| package. This
package can be downloaded in source from the \url{http://tug.org/}
website and compiled locally, or installed directly using a package
management command. In either case you will need to the installation
as the superuser; this is typically done with the |sudo| command.

For Debian and Ubuntu Linux, use the |apt-get| command:

\begin{Verbatim}
% sudo apt-get install texlive-latex3
\end{Verbatim}

For Fedora, RedHat, and Centos, use the |yum| command:

\begin{Verbatim}
% sudo yum install texlive-latex3
\end{Verbatim}

We have noticed that the install command occasionally fails. If it
fails for you, try it again two or three more times. If that still does not work, you will need to download and
install from source.

\subsection{Installation on Microsoft Windows}
One of the most popular distributions for Microsoft Windows machines is called
MiK\TeX, which can be downloaded from \url{http://www.miktex.org/}.
In most cases, this distribution of \LaTeX{} will automatically
download additional packages if needed.

\subsection{Creating classified documents}
If you are producing a classified thesis, you should ask a
system administrator at the STBL or SCIF to provide you with a
\LaTeX{} installation on a computer authorized to handle classified
information. With appropriate approvals you can alternatively set up an
installation on a stand-alone machine.
In the event that you are missing a package to create your thesis, the |.sty| files can be copied from the package on the CTAN website into your thesis directory.  A system administrator should be consulted on the file transfer process.  

\section{Additional Applications for \LaTeX{} }

In addition to downloading and installing \LaTeX, you will need an
editor for editing the |.tex| input files. These programs are
sometimes called \emph{front ends}, although the term is imprecise and
probably incorrect. You can use \emph{any} editor for editing a |.tex|
file, even Microsoft Word. Indeed, most of the front ends for \LaTeX{} 
are really just text editors with syntax highlighting, although some
will automatically compile your document and even jump to errors in
the input file when they are encountered.

\subsection{Selection of an Editor}
Text editors offer a variety of features.  Some are easier for beginners,
 such as Notepad++, LEd, and TeXnicCenter.  Advanced editors for skilled users
 include emacs, vim, and Texlipse.  A thorough comparison of editors 
 is available at
 \url{http://en.wikipedia.org/wiki/Comparison_of_TeX_editors}.  Learning about the various
 features of the text editors can help you to dramatically improve your writing efficiency.

If you are running \LaTeX{} on MacOS, you already have a powerful editor
installed on your computer: EMACS, which can be run from within the
Macintosh Terminal application. You can run the tutorial for EMACS by
starting Terminal, typing |emacs| and enter at the command prompt, and
then typing control-h followed by a ``|t|'' to start the tutorial.

\subsection{LyX: An alternative \LaTeX{} system}

LyX is a program that provides a WYSIWYG (What You See Is What
  You Get) graphical user interface for \LaTeX. Instead of editing the
|.tex| file directly, however, you edit an intermediate form which is
a restricted set of \LaTeX. LyX then runs \LaTeX{} for you and
produces the results.

While LyX is easier to use than \LaTeX, it does not have the power or
the flexibility.

LyX is free and open source and actively supported. LyX provides a GUI interface for floating figures and tables,
formatting, fonts, labels, chapters, sections, subsections, math
equations, tables, and much more.  It also has built-in features for
supporting Bib\TeX{} bibliographies, citations, and cross-references,
and generally supports anything that can be done with \LaTeX{}.

LyX can be downloaded from  \url{http://www.lyx.org/} and is available
for Linux, MacOS and Windows.  If you want spell-checking, you will
also need to install the \emph{aspell} package. 

An NPS thesis and dissertation template developed by CDR Michael
Bilzor is available for use with LyX. It can be be obtained from
\url{http://simson.net/npsthesis/lyxthesis.zip}.  

Further information about the LyX template can be found in the Appendix.


\section{Running \LaTeX}\label{runninglatex}
The \LaTeX{} system is actually a set of programs. For creating a thesis at NPS you will use several programs:

\begin{description}
\item{\texttt{pdflatex}} This program reads the input file (\eg
  |thesis.tex|) and produces a |.pdf| file (\eg |thesis.pdf|)
  as an output. This program also produces a number of intermediate
  files (|thesis.aux|, |thesis.bbl|, |thesis.toc|,
  \etc).
\item{\texttt{xelatex}} is a version of \LaTeX designed to process
  Unicode used in non-Roman languages. In some cases packages that are
  designed to work with \LaTeX will not work with \texttt{xelatex},
  which is why we do not recommend using it unless you have no other choice.
\item{\texttt{latex}} This is an older version of the |pdflatex|
  program that produces |.dvi| files. The |.dvi| file must then be
  transformed into either a |.ps| or a |.pdf| file. In practice you
  should not run |latex| unless you need to use a special graphics
  module called |PStricks|. That module is beyond the scope of this document.
\item{\texttt{bibtex}} This program reads the |thesis.bbl| file and
  produces a bibliography in a file called |thesis.bst| which
  includes the bibliography. The |thesis.bst| then gets read the
  next time |pdflatex| is run.
\item{\texttt{authorindex.pl}} This is a program in perl that produces the
  author index from the |thesis.bbl| file. The authorindex is
  saved in the file |thesis.ain|.
\item{\texttt{fixerrors.py}} It turns out that there is a bug in Bib\TeX{}
  which causes URLs longer than 53 characters to be improperly
  split. This program unsplits them. It also will correct |authorindex| errors.  You do not need to use this
  program if you do not have these errors.
\end{description}

If you are processing a file \emph{thesis.tex} to create a
\emph{thesis.pdf} file, you will typically execute these commands in
this order:

\begin{enumerate}
\item |pdflatex thesis|
\item |bibtex thesis|
\item |python fixerrors.py thesis|
\item |pdflatex thesis|
\item |pdflatex thesis|
\end{enumerate}

The first run of |pdflatex| creates the file |thesis.aux| (and any
other |.aux| files that might be needed).  A PDF file is also created,
but if you have any backwards references in your document the PDF file
will contain incorrect references. The subsequent runs read
these |.aux| files and generate correct back-references. Some editors, such as {\TeX}lipse, will run these commands as needed on behalf of the user.  Each command has a |thesis| argument, which is
 referring to the file |thesis.tex| without the |.tex| file extension.
 In rare cases it may be necessary to delete the |.aux| files and
 re-run the |pdflatex| command from the beginning.

\section{Basic \LaTeX{} formatting}
Here is a simple \LaTeX{} document:
\begin{Verbatim}
\documentclass{article}
\begin{document}
Hello World!
\end{document}
\end{Verbatim}

Normally with \LaTeX{}, you just type text and leave a blank line between
each paragraph. \LaTeX{} then formats it into beautiful
paragraphs. \LaTeX{} will ignore the space at the beginning of each line.

Here is a slightly more complex document:

\begin{Verbatim}
\documentclass{article}
\begin{document}
In December 1951, in a move virtually unparalleled in the history of
academe, the Postgraduate School moved lock, stock and wind tunnel
across the nation, establishing its current campus in Monterey,
Calif. 
     The coast-to-coast move involved 500 students, about 100
     faculty and staff and thousands of pounds of books and 
     research equipment. Rear Adm.\ Ernest Edward Herrmann 
     supervised the move that  pumped new vitality into the
     Navy's efforts to advance naval science and technology.

% This is comment. Nobody will see it.

Today the school, known as the ``Naval Postgraduate School,'' is the
Navy's preeminent institution of graduate education and advanced
research. Approximately 1 in 10 of the students are in the top 10\% of
their classes.
\end{document}
\end{Verbatim}


\LaTeX{} will format the above text into a document that looks  like this:

\fbox{
\begin{minipage}{6in}
\setlength{\parindent}{2pc}
In December 1951, in a move virtually unparalleled in the history of
academe, the Postgraduate School moved lock, stock and wind tunnel
across the nation, establishing its current campus in Monterey,
Calif. The coast-to-coast move involved 500 students, about 100
faculty and staff and thousands of pounds of books and research
equipment. Rear Adm.\ Ernest Edward Herrmann supervised the move that
pumped new vitality into the Navy's efforts to advance naval science
and technology.\\
\\
\indent Today the school, known as the ``Naval Postgraduate School,'' is the
Navy's preeminent institution of graduate education and advanced research.
Approximately 1 in 10 of the students are in the top 10\% of
their classes.
\end{minipage}
}

This sample document illustrates a few important points about \LaTeX:
\begin{itemize}
  \item \LaTeX{} will automatically
  re-wrap your text as necessary to format the paragraphs. Indention is determined by the style of the current document, not by whether or not you actually indent the paragraph.
\item \LaTeX{} ignores space at the beginning of lines; breaks between
  paragraphs are marked with blank lines.
\item What you type is not what you get!  In particular, opening
  double quote marks are typed as two backquotes (\verb+``+) and closing
  double quote marks are entered as two apostrophes (\verb+''+). You will
  also note that the period following Rear Adm.\ Ernest Edward
  Herrmann's name is followed by a backslash (\verb+\+) and a space, rather
  than just a space. This tells \LaTeX{} that the period does not mark
  the end of a sentence.
\item Commands begin with a backslash (|\|) and contain only uppercase
  and lowercase letters.
\item A command can have zero or more arguments. The arguments are
  enclosed within braces (|{}|). The |\documentclass| command begins
  the document; its argument is the kind of document you are making. 
\item The |\begin| command introduces an \emph{environment}. Every
  document has at last one environment, the |document|
  environment. Every |\begin| must have a matching |\end| that names
  the same environment. Environments can be nested.
\item Comments can be embedded in your document with a percent sign
  (|%|). Anything after the percent sign will not print. To print a
  percent sign, prefix it with a backslash (|\%|).
\end{itemize}

\section{Typing Special Characters}
This section provides information on how to type special characters in
\LaTeX. In each section we will have a table that shows what to type
and how it displays in your final |.pdf| file.  If the above text were
put into a table, it would look like this:

\begin{center}
\begin{tabular}{l|l||l|l}
Typed & Displayed & Typed & Displayed \\\hline
|``|  & ``        & |''|  & '' \\
\end{tabular}
\end{center}

The left entry of the table is the backquotes, which shares the tilde key on the US keyboard.
The right entry is the single quotes, which shares the double quote marks key.

\subsection{Typing Quotes}

To type quotes, you should not use the double-quote character. Instead,
use the back quote (|`|) and the forward quote (|'|)
to type quotes:

\begin{center}
\begin{tabular}{l|l||l|l}
Typed               & Displayed & Typed   & Displayed\\\hline
|don't|             & don't     &  |3'2''|  & 3'2''\\
|``this''|          & ``this''  & |`is'|  & `is' \\
|``\,`special'\,''| & ``\,`special'\,'' \\
\end{tabular}
\end{center}


\subsection{Controlled Special Characters}\label{special}

Unlike Microsoft Word and other programs, \LaTeX{} uses special
characters embedded in your text to control formatting. The most
common of these characters is the backslash ($\backslash$). All of
special characters are listed below:

\begin{center}
\begin{tabular}{c|p{4in}}
Special Character & Why it is special\\
\hline
|\|      & Introduces a command\\
\verb|{| & Introduces arguments in commands or the start of a group\\
\verb|}| & Closes arguments in commands or the end of a group\\
|%|      & The comment character; \LaTeX{} ignores the rest of the line\\
|#|      & Used for parameter substitution inside macros\\
|~|      & Enters a hard, non-breaking space\\
|&|      & Used for delimiting columns in a table \\
\verb|$| & Turns on/off math mode (see \S\ref{math})\\
|_|      & Used for subscript in math mode\\
|^|      & Used for superscript in math mode\\
\end{tabular}
\end{center}

To enter the special characters into your
document you must use a special sequence that begins with a
backslash. Most (but not all) of these special sequences are the
character itself. If you are curious, inside \LaTeX{}, each of these
sequences is implemented as a command that causes \LaTeX{} to output the
character that has been quoted:

\begin{center}
\begin{tabular}{l|l||l|l}
Typed  & Displayed & Typed          & Displayed\\\hline
|\$|   & \$        & |\&|           &    \&   \\
|\{|   & \{        & |\}|           &    \}   \\
|\%|   & \%        & |\_|           &    \_   \\
|\#|   & \#        & |\^{}|         &    \^{} \\
|\~{}| & \~{}      & |$\backslash$| & $\backslash$ \\
\end{tabular}
\end{center}


\subsection{Accented, Dotless and Slashed Vowels}
With \LaTeX{} most accented vowls are entered with a combining accent character and a
vowel, although some (such as the angstrom symbol) are not, as shown below:

\begin{center}
\begin{tabular}{l|l||l|l}
Typed    & Displayed & Typed   & Displayed\\\hline
|\'{o}|  & \'{o}     & |\~{o}| & \~{o} \\
|\'{o}|  & \'{o}     & |\={o}| & \={o} \\
|\^{o}|  & \^{o}     & |\.{o}| & \.{o} \\
|\"{o}|  & \"{o}     & |\d{o}| & \d{o} \\
|\c{o}|  & \c{o}     & |\u{o}| & \u{o} \\
|\b{o}|  & \b{o} \\
|\aa|    & \aa       & |\AA|   & \AA \\
|\i|     & \i        & |\j|    & \j \\
|\o|     & \o        & |\O|    & \O \\
|\ae|    & \ae       & |\AE|   & \AE \\
|\oe|    & \oe       & |\OE|   & \OE \\
|\v{o}|  & \v{o}     & |\H{o}| & \H{o} \\
|\t{oo}| & \t{oo} \\
\end{tabular}
\end{center}

Note that |\i| displays a dotless i while |\j| displays a dotless
j. Some fonts do not have some of these characters, and display a
black box instead.

\subsection{Symbols}
\LaTeX{} has literally hundreds of symbols that you can include in
your document. These symbols  include the copyright symbol, currency symbols, foreign language characters, and many more.  
The symbols are placed in documents using macros, allowing 
a plain text document to support a large variety of non-standard text characters.  A
complete guide to the available symbols in \LaTeX{} is available online at 
\url{http://mirror.ctan.org/info/symbols/comprehensive/symbols-letter.pdf} .

Here are some of the most common symbols you are likely to need for an
NPS document:

\begin{center}
\begin{tabular}{l|l||l|l}
Typed & Displayed & Typed & Displayed\\\hline
|\l|     & \l    & |\L|         & \L \\
|\S|     & \S    & |\ss|        & \ss \\
|\P|     & \P    & |\pounds|    & \pounds \\
|?`|     & ?`    & |\copyright| & \copyright \\
|!`|     & !`    & |\texttrademark| & \texttrademark\\
|\euro|  & \euro & |\textregistered| & \textregistered\\
|\dag|   & \dag  & |\ddag| & \ddag \\
\end{tabular}
\end{center}

(Note: The |\euro| command requires that the command
|\usepackage{eurosym}| be part of the document's \emph{preamble} (the
part before the |\begin{document}|). It is included as part of the
|npsreport.cls| file.)

\section{Fonts}
Like Microsoft Word, \LaTeX{} makes it easy to alter font, size, face,
and weight of text. But unlike Word, these changes are typically done
in a structured manner that lends itself to creating documents that
have consistent font usage throughout. 

\subsection{Changing Font Size}
Although \LaTeX{} allows you to use fonts of any size, the built-in
templates provides eleven built-in sizes. These sizes automatically
adjust depending on if you are creating a document with 12-point font
(NPS standard), 11-point font, or 10-point font. 

When you change the font size, that change stays in effect until you
change it again. You can confine your font change by placing the text
you want resized within braces, sometimes called a \emph{group} or a
\emph{block}, as shown in the examples in the
following table:

\begin{center}
\begin{tabular}{llll}
Size          & Point Size & Typed & Displayed  \\\hline
|\tiny|       &  6 & |{\tiny This is tiny}| & {\tiny This is tiny}\\
|\scriptsize| &  8 & |{\scriptsize scriptsize}| & {\scriptsize scriptsize} \\\
|\footnotesize| & 10 & |{\footnotesize footnotesize}| & {\footnotesize footnotesize} \\\
|\small|        & 11 & |{\small This is small}| & {\small This is small} \\\\\
|\normalsize|   & 12 & |{\normalsize normalsize}| & {\normalsize  normalsize} \\\\\\\
|\large|        & 14 & |{\large This is large}| & {\large This is large} \\\\\\\
|\Large|        & 17 & |{\Large Large}| & {\Large Large} \\\\\\\
|\LARGE|        & 20 & |{\LARGE LARGE}| & {\LARGE LARGE} \\\\\\\
|\huge|         & 25 & |{\huge huge}| & {\huge huge} \\\\\\\
|\Huge|         & 25 & |{\Huge Huge}| & {\Huge Huge} \\\\\\\
\end{tabular}
\end{center}

Notice that with the NPS template there is no difference between
|\huge| and |\Huge|.

You can also pick an arbitrary size by using the |\fontsize| and
|\selectfont| commands. The |\fontsize| command takes two arguments:
the size of the font and the size of the leading, or the amount of
space between lines. First the size is selected with the
|\fontsize{i}{j}| command where |i| and |j| are expressed in points
(there are 72.27 points in an inch). Next the font is selected with the |\selectfont| command, as shown below:

\begin{center}
\begin{tabular}{l|l}
Typed & Displayed  \\\hline
|{\fontsize{4}{5}\selectfont very tiny}|& {\fontsize{4}{5}\selectfont very tiny}\\
|{\fontsize{64}{64}\selectfont Big}|& {\fontsize{64}{64}\selectfont Big}\\
\end{tabular}
\end{center}

\subsection{Changing Font Style}
\LaTeX{} provides these macros for selecting font styles:

\begin{center}
\begin{tabular}{l|l}
Typed                                  & Displayed\\\hline
|\textrm{This is Roman}|               & \textrm{This is Roman} \\
|\textbf{This is bold}|                & \textbf{This is bold}\\
|\texttt{This is typewriter}|          & \texttt{This is typewriter}\\
|\textsc{This is small capitals}|      & \textsc{This is small capitals}\\
|\textsl{This is slanted}|             & \textsl{This is slanted}\\
|\textsf{This is sans serif}|          & \textsf{This is sans serif}\\
|\textit{Italics}|                     & \textit{Italics}\\
|\emph{This is emphasized}|            & \emph{This is emphasized}\\
|$\cal CALLIGRAPHICS$|                 & $\cal CALLIGRAPHICS$\\
|{\boldmath $\cal BOLD CALIGRAPHICS$}| & {\boldmath $\cal BOLD CALIGRAPHICS$}\\
\end{tabular}
\end{center}

Notice that the last two are surrounded by dollar signs as they require math mode (see \S\ref{math}).

If you just want to put something into italics, you should use
|\emph{text}|, which will produce \emph{text}. The reason to use
|\emph{}| and not |\textit{}| is that |\emph{}| will nest as
necessary.  For example, this:

\begin{center}
\begin{Verbatim}
\emph{You can even \emph{emphasize} a word within a sentence
that is emphasized.}
\end{Verbatim}
\end{center}

typesets as this:

\begin{center}
\emph{You can even \emph{emphasize} a word within a sentence that is emphasized.}
\end{center}

As you should in general avoid underlining text, we will not show you
how to do it in this document.

\subsection{Choosing an Arbitrary Font}
There are many ways that you can request arbitrary fonts for small
sections of your document, but they are all beyond the scope of this
article. It is also possible to embed arbitrary Unicode within a
\LaTeX{} document, either by using the Unicode-aware version of
\LaTeX{} called \raisebox{-3pt}{\includegraphics[height=12pt]{xelatex}} (|xelatex|), or by saving your Unicode characters in a |.pdf| file and
including that file with the |\includegraphics{}| command (as we
did with the \raisebox{-3pt}{\includegraphics[height=12pt]{xelatex}} logo).

If you want to change the font for an entire document, please refer to the
\LaTeX{} Font Catalogue at \url{http://www.tug.dk/FontCatalogue/}
which provides documentation and examples.  To change the font of the
NPS template, refer to \S\ref{thesisprologue}.

\section{Math}\label{math}

Typsetting mathematics is one of the primary design goals of
\LaTeX. The program  has more features for typsetting math
than typsetting text. There is also a powerful set of mathematical
extensions by the American Mathematical Society called |amsmath|. 
For more detailed information, please see:

\begin{itemize}
\item The \LaTeX{} wikibook,
  \url{http://en.wikibooks.org/wiki/LaTeX/Mathematics}.  
\item The User's guide for the \texttt{amsmath} package is available at:
\url{ftp://ftp.ams.org/ams/doc/amsmath/amsldoc.pdf}
\item The short math guide, available at
  \url{http://tinyurl.com/63w3mnu}.
\end{itemize}

What follows here is necessarily very brief. 

\subsection{Math Mode}
To typset math you must enter math mode. There are two easy ways to
enter math mode. 

\begin{itemize}
\item You can put your math between two dollar signs. For example,
  entering |$1+1=2$| in your document will produce $1+1=2$.
\item You can put the equation on a line by itself in an \emph{equation}
  environment (\eg between |\begin{equation}| and |\end{equation}|
  commands). An equation environment creates a block that is typeset in
  math mode and includes a numbered equation. For example, this:
\begin{Verbatim}
\begin{equation}
1+1=2
\end{equation}
\end{Verbatim}

produces this:

\begin{equation}
1+1=2
\end{equation}
\end{itemize}

Math mode can also be used in paragraphs to add special math characters, such as the $\pi$ symbol (using |$\pi$| here).
In fact, many of the symbols that \LaTeX{} displays can only be
displayed while in math mode.

\subsection{Simple Math in Math Mode}
As the examples above show, you can type simple math in
math mode and get what you want. In general, variables (the letters a
through z), the plus (|+|), minus (|-|) and equals (|=|) all 
typeset properly when you type them between dollar signs.  But there are some caveats:

\begin{itemize}
\item If you wish to typeset a multiplication symbol, use |\times| instead
of an asterisk (|*|).
\item If you wish to typeset division, use |\div| to enter a division
  symbol or the |\frac| command to create a fraction. Do not use |/|
  for division. 
\item Spaces are ignored in math mode. If you want a space, you
  probably should use multiple equations, with each equation
  in math mode but with a non-math mode space between them.
\end{itemize}

The table below shows some examples:

\begin{center}
\begin{tabular}{l|l|l}
Typed                   & Displayed & Comments \\\hline
|$1+2=3$|               & $1+2=3$ \\
|$10 * 10 = 100$|       & $10 * 10 = 100$  & {\small Don't use asterisks for multiplication.}\\
|$10\times10=100$|      & $10\times10=100$ & Spaces don't matter\\
|$10  \times  10  =  100$|      & $10  \times  10  =  100$\\
|$a=3$|                 & $a=3$\\
|$a=f/m$|               & $a=f/m$ & {\small Don't use the slash for division.}\\
|$a=f \div m$|          & $a=f \div m$\\
|$a=\frac{f}{m}$|       & $a=\frac{f}{m}$\\
|$f(x)=3x$|             & $f(x)=3x$\\
|$a=2b$ and $a+2=b+4$|  & \multicolumn{2}{l}{$a=2b$ and $a+2=b+4$}\\
\end{tabular}
\end{center}

Use math mode when you need to enter math---it's worth the effort. For example, consider a function of t. You can certainly type this
without math mode---witness f(t)---but doesn't it look much better
when dollar signs are placed around the symbol, like this: $f(t)$? Math
mode adds clarity.

If you wish to discuss an important equation in your document,
use the \emph{equation}
environment.  This environment sets your equation off from the body
text and gives it a number.

Using the equation environment, this:

\begin{Verbatim}
\begin{equation}
a = 1+2
\end{equation}
\end{Verbatim}

typesets as:

\begin{equation}
a = 1+2
\end{equation}



\subsection{Superscripts and subscripts}
In math mode the caret (|^|) is used for superscript and the
underbar (|_|) is used for subscript (this is why the characters are
special). The commands only superscript or subscript the following
character; if you want to superscript or subscript multiple characters
you need to make them a group by enclosing them in braces.

Here are some examples:

\begin{center}
\begin{tabular}{l|l}
Typed                   & Displayed\\\hline
|$a^2+b^2=c^2$|         & $a^2+b^2=c^2$\\
|$2^{16}=65,535$|       & $2^{16}=65,535$\\
|$N_A$ is Avogadro's constant| & $N_A$ is Avogadro's constant \\
|$A^{B^C}$|             & $A^{B^C}$\\
|$A^{B^{C^D}}$|             & $A^{B^{C^D}}$\\
|$a_k$ and $b_k$|       & $a_k$ and $b_k$         \\
|$a_k$ and $b_k$|               & $a_k$ and $b_k$ \\
\end{tabular}
\end{center}

\subsection{Combining Symbols in Groups}

Many math symbols use subscripts and superscripts to determine
placement of specific equation elements.  This includes |\int| which is used to create integrals
and |\sum| which is used to create sums. Below are some examples with the Fourier series.

\begin{center}
\begin{tabular}{l|l}
Typed                   & Displayed\\\hline
\\
\verb|$a_k=\frac{1}{\pi}|                  & $a_k = \frac{1}{\pi}\int_{-\pi}^{\pi}f(x)cos(kx)dx$\\
\verb|\int_{-\pi}^{\pi}f(x)cos(kx)dx$|      & \\[4pt]
\hline
\\
\verb|$b_k=\frac{1}{\pi}|                  & $b_k = \frac{1}{\pi}\int^{\pi}_{-\pi}f(x)\sin(kx)dx$\\
\verb|\int^{\pi}_{-\pi}f(x)\sin(kx)dx$|      & \\[4pt]
\hline
\\
\verb|$f(x)=\sum_{k=-\infty}^{\infty}| & $f(x)=\sum_{k=-\infty}^{\infty}c_k e^{-jkx}$\\
\verb|c_k e^{-jkx}$|                 & \\[4pt]
\end{tabular}
\end{center}

The |\sum| and |\int| symbols display these elements differently in an \emph{equation} environment:

\begin{equation}
f(x)=\frac{a_0}{2} + \sum_{n=1}^{\infty}\left[a_k \cos(kx) + b_k \sin(kx)\right]
\label{fourier}
\end{equation}

We used this code to type \eqnref{fourier}:
\begin{Verbatim}
\begin{equation}
f(x)=\frac{a_0}{2} + \sum_{n=1}^{\infty}
    \left[a_k \cos(kx) + b_k \sin(kx)\right]
\label{fourier}
\end{equation}
\end{Verbatim}

There are several important math mode conclusions to draw from these examples:
\begin{itemize}
  \item Simple one letter subscripts and superscripts do not need to be enclosed by curly braces, but multiple character ones must.
  \item Subscripts and superscripts in math mode are to the right of
    the sum and integral characters; however, in the equation
    environment they are above and below.  In the other uses
    (exponentials, \etc), the results are the same in either mode.
  \item The order of appearance of subscript than superscript or vice versa does not change the results.
\end{itemize}


The |amsmath| package has additional environments, symbols and
commands such as provisions for non-numbered,
multiple-lined and aligned equations. Its user guide is an
excellent reference and provides many examples.

Although math mode may seem cumbersome at first, its syntax does become second nature and very sophisticated equations can be generated, if needed.
Consult the user manuals and references provided in \S\ref{math}.

\subsection{Parenthesis and Brackets}
You can use regular parentheses in math mode, but they do not stack nicely:

\begin{center}
\begin{tabular}{l|l}
Typed                   & Displayed\\\hline
|$((1))$| & $((1))$\\
|$[((1))]$| & $[((1))]$
\end{tabular}
\end{center}


You probably want the outer brackets and parenthesis to be bigger than
the inner ones. You can do that using the |\left| and |\right|
commands.  Following these macros is another character, such as a
parenthesis or bracket.  
These macros place the correctly sized
specified character into the equation.  They will automatically get
bigger as necessary, especially in the \emph{equation} environment.

For example, here is what you get with the conventional parenthesis:
\begin{Verbatim}
\begin{equation}
A = ( \sum_{i=1}^{10}i\times\sin(i) )
\end{equation}
\end{Verbatim}

\begin{equation}
A = ( \sum_{i=1}^{10}i\times\sin(i) )
\end{equation}

And here is an example using |\left| and |\right|:

\begin{Verbatim}
\begin{equation}
A = \left( \sum_{i=1}^{10}i\times\sin(i) \right)
\end{equation}
\end{Verbatim}

\begin{equation}
A = \left( \sum_{i=1}^{10}i\times\sin(i) \right)
\end{equation}

This example  seems silly, but shows  how powerful
|\left| and |\right|  are:

\begin{Verbatim}
\begin{equation}
\left(\frac{
    \left(\frac{1}{2}\right)}{
    \left(\frac{3}{4}\right)
}\right)
\end{equation}
\end{Verbatim}

Produces this:

\begin{equation}
\left(\frac{
    \left(\frac{1}{2}\right)}{
    \left(\frac{3}{4}\right)
}\right)
\end{equation}


%%%%%%%%%%%%%%%%%%%%%%%%%%%%%%%%%%%%%%%%%%%%%%%%%%%%%%%%%%%%%%%%

\section{Spacing, Frameboxes, and Centering}
\LaTeX{} has a number of commands for controlling space, creating
boxes, and centering text.

\subsection{Controlling Spaces}

You can use these commands for controlling how much space is inserted
between words:

\begin{center}
\begin{tabular}{l|l}
Typed                                 & Displayed\\\hline
|\,| (a slash followed by a comma)    & produces\,a\,small\,space. \\
|\ | (a slash followed by a space)    & produces\ a\ standard\ word\ space. \\
|\@| (a slash followed by an at sign) & produces an intersentence space. \\
\end{tabular}
\end{center}


The small space and standard word space are used between the words in the right column.  An intersentence space is needed to correct the following error where a sentence ends in a capital letter.

\subsection{Suppressing Orphans and Widows}
When typesetting, paragraphs that have
their last line on the following page are called \emph{orphans}, and
paragraphs that begin at the bottom of a page with a single line are
called \emph{widows}. Orphans and widows are considered
ugly. Microsoft Word can be programmed to prevent widows and orphans
by requiring that all paragraphs have at least 2, 3 or 4 lines on a
page. Microsoft Word achieves this desired result by inserting spaces
between paragraphs to balance out the page as necessary. \LaTeX{} has
no automatic control over orphans and widows. It is one of the major
failings of the system.

NPS students working on their Master's thesis are frequently
instructed by the Thesis Processor to modify their document so that
there are no widows or orphans. One way to achieve such a
result is to rewrite the text by inserting sentences or removing
them. This may seem excessive to you. Another way you can control
widows and orphans is by manually adding or removing space, or by
stretching a page to allow additional lines on it.

You can force a blank line using |\\| and have the option to force a
specific length using |\\[3pt]|:
\begin{center}
\begin{tabular}{l|l}
Typed        & Displayed\\\hline
|This is\\|  & This is\\
|an example.| & an example.\\
\\
|This is\\[3pt]|  & This is\\[3pt]
|an example.| & an example.
\end{tabular}
\end{center}

You can add or remove space on a page with the |\enlargethispage| command. For
example, to squeeze another line onto the current page, insert this
command onto the page:

\begin{verbatim}
    \enlargethispage{1pc}
\end{verbatim}

This command is useful when you need to manually enlarge a
page so that the last line of a paragraph can fit on the present page
without being pushed to the next page. 

You can shorten a page by a line, forcing a widow onto the next page:

\begin{verbatim}
    \enlargethispage{-1pc}
\end{verbatim}


These commands are also useful when writing a conference paper that
needs to fit within a certain page length.

\subsection{Frameboxes and Centering}

You can draw a box around text with:

|\framebox[width]{textstring}|

The |[width]| parameter is optional. Without it, the box defaults to
the minimum size necessary to hold the \texttt{textstring}.

Here are is an example to show what we mean. This text:

\begin{Verbatim}
\framebox{This is an important statement.}\\
\framebox[15pc]{This is an important statement.}\\
\framebox[30pc]{This is an important statement.}
\end{Verbatim}

will typeset as this:

\framebox{This is an important statement.}\\
\framebox[15pc]{This is an important statement.}\\
\framebox[30pc]{This is an important statement.}

If you provide a space that is too small, the results will be ugly:

\begin{Verbatim}
\framebox[5pc]{This is an important statement.}
\end{Verbatim}

\framebox[5pc]{This is an important statement.}

You can center any text, table, figure, \etc with the \emph{center} environment, using the following:

\begin{Verbatim}
\begin{center}
\framebox{$a^2 + b^2 = c^2$}
\end{center}
\end{Verbatim}

Produces this:

\begin{center}
\framebox{$a^2 + b^2 = c^2$}
\end{center}

If you are entering a number of equations in your document, you may
want to use the \texttt{equation} environment, which will provide
numbered equations:

\begin{Verbatim}
\begin{equation}
a^2 + b^2 = c^2
\end{equation}
\end{Verbatim}
\begin{equation}
a^2 + b^2 = c^2
\end{equation}

Unlike many things in \LaTeX, the \emph{framebox} and \emph{equation}
environments cannot be readily combined. That's because |\fbox| turns
off math mode, so you need to manually turn it back on.

Type this:

\begin{Verbatim}
\begin{equation}
\fbox{$a^2 + b^2 = c^2$}
\end{equation}
\end{Verbatim}


To produce this:

\begin{equation}
\fbox{$a^2 + b^2 = c^2$}
\end{equation}



\section{Lists}

There are three kinds of lists that you may wish to make:

\begin{description}
\item[description] lists are used for definitions, where a  short phrase is bolded and the remainder text is the standard font (like this list).
 \item[enumerate] lists are lists where each item is numbered and the
  ordering is relevant, like the steps of a recipe. 
\item[itemize] lists are lists where each item is of equal
  importance.
\end{description}

Lists are implemented as \LaTeX{} environments, which means that they
begin with a |\begin{|\emph{listname}|}| and end with an
|\end{|\emph{listname}|}|.

\begin{minipage}{.45\textwidth}
\begin{Verbatim}
\begin{description}
\item[Earth] Third Planet.
\item[Mars] Fourth Planet.
\item[Venus] Second Planet.
\end{description}
\end{Verbatim}
\end{minipage}
\hfill
\begin{minipage}{.45\textwidth}
\begin{description}
\item[Earth] Third Planet.
\item[Mars] Fourth Planet.
\item[Venus] Second Planet.
\end{description}
\end{minipage}


\begin{minipage}{.45\textwidth}
\begin{Verbatim}
\begin{enumerate}
\item Wake up.
\item Go to work.
\item Go home.
\item Go to sleep.
\item Repeat.
\end{enumerate}
\end{Verbatim}
\end{minipage}
\hfill
\begin{minipage}{.45\textwidth}
\begin{enumerate}
\item Wake up.
\item Go to work.
\item Go home.
\item Go to sleep.
\item Repeat.
\end{enumerate}
\end{minipage}

\begin{minipage}{.45\textwidth}
\begin{Verbatim}
\begin{itemize}
\item Hamburgers
\item Hotdogs
\item Chips
\end{itemize}
\end{Verbatim}
\end{minipage}
\hfill
\begin{minipage}{.45\textwidth}
\begin{itemize}
\item Hamburgers
\item Hotdogs
\item Chips
\end{itemize}
\end{minipage}


\section{Labels and Captions}\label{sec:labels}
Labels are hidden markers in your |.tex| files created by the
|\label{name}| command.  These markers are never shown directly in the
output files.  However, correctly placing these markers in your file
allows you to reference chapters, appendices, sections, figures,
tables and equations.  You may wish to give your marker names prefixes
such as \texttt{chap:}, \texttt{sec:}, \texttt{fig:}, \texttt{tab:}
and \texttt{eqn:} to logically identify the labels.

Captions are the text that appears below a figure or table to provide
context for the information presented.  Captions are indicated by the
command \verb|\caption{Sentence.}| Sometimes longer captions can look
poorly in the Lists of Figures and Tables, so a caption for the table
can also be specified by using
\verb|\caption[Short Sentence.]{Longer Sentence.}|

\section{Tables}
Tables provide a valuable means to display data in an organized
manner.  Using tables in \LaTeX{} is easy once the syntax is
understood, although tables can be incredibly complex as well.  To see
plenty of examples and explanation of options, visit
\url{http://en.wikibooks.org/wiki/LaTeX/Tables}. There are also
\emph{many} add-on packages which provide additional functionality.

Here is a very simple table:

\begin{Verbatim}
\begin{center}
\begin{tabular}{cc}
Shape & Sides \\
\hline
Triangle & 3 \\
Square & 4 \\
\end{tabular}
\end{center}
\end{Verbatim}

Produces this:

\begin{center}
\begin{tabular}{cc}
Shape & Sides \\
\hline
Triangle & 3 \\
Square & 4 \\
\end{tabular}
\end{center}

This table has six things we have not seen before:

\begin{enumerate}
\item The |{cc}| is the table format specifier. Each ``c'' represents
  a column that is centered.
\item The |&| is used to separate columns of the table.
\item The |\\| is used to separate rows of the table.
\item The |\hline| is used to draw a horizontal line.
\item The table is drawn in a |tabular| environment.
\item The |center| environment is used to center the table.
\end{enumerate}

Here is a slightly more complex table:

\begin{Verbatim}
\begin{center}
\begin{tabular}{lcr}
Left Justified & Center Justified & Right Justified \\
\hline
Each & Column & Entry \\
Spaces      & Do not    & Matter \\
But         & May       & Assist you. \\
\end{tabular}
\end{center}
\end{Verbatim}

And here is how it formats:

\begin{center}
\begin{tabular}{|lcr}
Left Justified & Center Justified & Right Justified \\
\hline
Each & Column & Entry \\
Spaces      & Do not    & Matter \\
But         & May       & Assist you. \\
\end{tabular}
\end{center}

This table has three columns as identified by the format specification
|lcr|. The first column is left-justified (``l''), the middle column
is center justified (``c'') and the third column is right justified
(``r''). The left side of the table has a vertical bar due to the pipe character (\verb+|+).  The |\hline| macro causes a horizontal line to be drawn across the table. 
Columns are indicated by the ampersands (|&|) and the amount of spaces used are not important.  Thus, the spaces can be used to ensure your columns line up in your |.tex| file the way the table is intended to print.
An empty cell is just ampersands separated by a space.
To start a new row of the table, use the double backslash (|\\|); this can be done at the end of the current row.

Below is a more complex table:


\begin{Verbatim}
\begin{center}
\begin{table}
\begin{tabular}[t]{|l||r|r|r|r|r|r|r|}
\hline
Pulse  & Pulse  & PRF & Duty  & Coupler & Power & Average & Peak  \\
Width  & Length &     & Cycle & Losses  & Meter & Power   & Power \\
\hline
Long & $1.0 * 10^{-6} s $ & 750 Hz & $7.0 * 10^{-4} $ & 
       49.1 dB & 3.4 dBm & 52.5 dBm & 73 dBm \\
& & & 31.075 dB & & & & 43 dBW \\
& & & & & & & 227.7 kW \\
\hline
Medium & $3.0 * 10^{-7} s $ & 1200 Hz & $2.0 * 10^{-4} $ & 
         49.1 dB & 0 dBm & 49.1 dBm & 83 dBm \\
& & & 34.29 dB & & & & 53 dBW \\
& & & & & & & 218.27 kW \\
\hline
Short & $1.0 * 10^{-7} s $ & 2400 Hz & $1.9 * 10^{-4} $ & 
        49.1 dB & -1 dBm & 48.1 dBm & 83 dBm \\
& & & 35.05 dB & & & & 53 dBW \\
& & & & & & & 206 kW \\
\hline
\end{tabular}
\end{center}
\end{Verbatim}

And here is the formatted table:

\begin{center}
\begin{tabular}[t]{|l||r|r|r|r|r|r|r|}
\hline
Pulse  & Pulse  & PRF & Duty  & Coupler & Power & Average & Peak  \\
Width  & Length &     & Cycle & Losses  & Meter & Power   & Power \\
\hline
Long & $1.0 * 10^{-6} s $ & 750 Hz & $7.0 * 10^{-4} $ &
       49.1 dB & 3.4 dBm & 52.5 dBm & 73 dBm \\
& & & 31.075 dB & & & & 43 dBW \\
& & & & & & & 227.7 kW \\
\hline
Medium & $3.0 * 10^{-7} s $ & 1200 Hz & $2.0 * 10^{-4} $ &
         49.1 dB & 0 dBm & 49.1 dBm & 83 dBm \\
& & & 34.29 dB & & & & 53 dBW \\
& & & & & & & 218.27 kW \\
\hline
Short & $1.0 * 10^{-7} s $ & 2400 Hz & $1.9 * 10^{-4} $ & 
        49.1 dB & -1 dBm & 48.1 dBm & 83 dBm \\
& & & 35.05 dB & & & & 53 dBW \\
& & & & & & & 206 kW \\
\hline
\end{tabular}
\end{center}

\section{Graphics}
This section briefly describes how to embed graphics in a \LaTeX{} document. For alternative
treatments we recommend H\"oppner's  ``Strategies for including graphics in \LaTeX{} documents''\cite{strategies}, and the book \emph{The LaTeX Graphics Companion}\cite{graphics-companion}.

Graphics are embedded using the |\includegraphics| command. It looks like this:

\begin{Verbatim}
       \includegraphics[options]{filename}
\end{Verbatim}

Typical options that you can use are |width=XX| or |height=XX|. For example, to include an image, have it centered, and scale it to 3 inches, use this:

\begin{Verbatim}
\begin{figure}
    \begin{center}
	\includegraphics[width=3in]{imagedirectory/file}
	\caption{Caption of the important figure.}
	\label{fig:importantfigure}
    \end{center}
\end{figure}
\end{Verbatim}

The |imagedirectory/file| notation indicates the images are in a subdirectory of |imagedirectory| and the file name is |file|.  The file extensions (|.jpg|, \etc) are optional. 

\emph{Note: \LaTeX{} does not properly include graphics that have
  periods in their filenames in addition to the period that is used to
  denote the file type!} 

There are two kinds graphics that you can include in a \LaTeX{}
document: 
\begin{description}
\item[Vector graphics] are excellent for printing purposes.  These images
 show up nicely on the computer screen and paper.  The computer renders these dynamically.  Several tools to generate vector graphics include |Inkscape| and
 \texttt{asymptote}. In general, you should convert vector graphics to a |.pdf| file and embed it in your document using |\includegraphics|.
\item[Raster graphics] include |.bmp|, |.jpg|, |.png|, and |.gif| images.  Formats without compression, including |.bmp| files
 generally will make your thesis file much larger.  Some of these images show up poorly or 
 pixelated in print copy.  You should convert these files to |.jpg| or |.png| and embed them in your document with |\includegraphics|.
\end{description}

Grahpics can also be created with R, gnuplot, matplotlib, or the
asymptote package.

Specific macros for including graphics provided by the |npsreport|
template are discussed in \secref{graphics}.

\section{Bib\TeX{} and Citations}
Bibliography and citation are important in your thesis.  Each department has different expectations
on citation styles.  Even if you select the wrong style initially, changing the format to the correct
option is only a one-line change for \LaTeX{}---unlike a complete nightmare for Microsoft Word users.

See \url{http://en.wikibooks.org/wiki/LaTeX/Bibliography_Management} and
 \\ \url{http://stefaanlippens.net/bibentry} for more information on BiB\TeX.
 
Many common citation reference organizers and websites output Bib\TeX{} formatted citations in |.bib| files.  This is a time saver,
as rather than retyping information, you can download the |.bib| file and copy its contents into your thesis master |.bib| file.
The master |.bib| file contains all of your citations, even ones that you have not yet cited.  When you reference these citations in your thesis, Bib\TeX{} can 
generate the references list for you.  To cite a specific reference use the \verb|\cite{name}| command.  After your first run of |pdflatex|, you follow with the |bibtex| command 
as in \S\ref{runninglatex}. 

\section{Going Further}
You are now on your way to becoming a \LaTeX{} expert and will find that many of the \LaTeX{} modules are already installed with your \LaTeX{}. 
You may find the following packages useful:

\begin{description}
\item[multirow] Allows a single table cell to extend to multiple
rows.

\item[ifthen]  Allows you to put conditions in your thesis. It is a
bit easier than using the |if| that's built in to \TeX.

\item[acronym]  A great package for automatically generating acronym lists.
It can track the usage of acronyms to correctly use the long name of the acronym
on its first use.

\item[makeidx]  For creating an index of key words and phrases in your document and what page the keywords appear on.

\item[mcode]  Places a MATLAB .m file contents directly into the output.  The
package performs syntax highlighting and line numbering if desired.  This package is not in the CTAN
and documentation may be obtained from the MATLAB Central website at \url{http://tinyurl.com/3wgcufr}.

\end{description}

In learning about \LaTeX{}, you are likely to do searches on the Internet, learn
about a new package, and want to learn more about it. You will
probably be surprised to find that many of these packages are already
part of your \LaTeX{} distribution. In most cases, the documentation for
the packages is already on your hard drive as a |.pdf| file.  We
recommend reading package documentation; the documentation frequently
has better information than the random pages you may find by using Google.

